% =========================================================================
% FUNDAMENTOS MATEMÁTICOS DOS GRAFOS DE FLUXO
% =========================================================================
\chapter{Fundamentos Matemáticos dos Grafos de Fluxo}\label{sec:base_matematica}

A formulação do problema de fluxo apoia-se na Teoria dos Grafos e na Otimização Combinatória (veja~\cite{ahuja1993network, bondy2008graph}). Uma rede de fluxo é um digrafo $D = (V, A)$, onde $V$ representa o conjunto de vértices e $A$ o conjunto de arcos.

A rede possui dois vértices terminais: uma fonte $s \in V$ e um sumidouro $t \in V$. Cada arco $(u,v) \in A$ possui uma capacidade não negativa $c(u,v) \ge 0$, dada por uma função $c: A \rightarrow \mathbb{R}_{\ge 0}$, que delimita o limite superior de fluxo admissível no arco.

A Figura~\ref{fig:rede_fluxo} ilustra uma rede de fluxo elementar com quatro vértices, onde as capacidades estão indicadas sobre os arcos.

\begin{figure}[!ht]
\centering
\begin{tikzpicture}[node distance=2.5cm, thick]
    \node[source node] (s) at (0, 0) {$s$};
    \node[flow node] (u) at (2.6, 1.5) {$u$};
    \node[flow node] (v) at (2.6, -1.5) {$v$};
    \node[sink node] (t) at (5.2, 0) {$t$};

    \draw[flow edge] (s) -- node[edge lbl above] {$10$} (u);
    \draw[flow edge] (s) -- node[edge lbl below] {$8$} (v);
    \draw[flow edge] (u) -- node[edge lbl, right] {$5$} (v);
    \draw[flow edge] (u) -- node[edge lbl above] {$7$} (t);
    \draw[flow edge] (v) -- node[edge lbl below] {$10$} (t);
\end{tikzpicture}
\caption{Rede de fluxo com fonte $s$, sumidouro $t$ e capacidades nos arcos.}
\label{fig:rede_fluxo}
\end{figure}

Um fluxo em $D$ é uma função $f: A \rightarrow \mathbb{R}_{\ge 0}$. O fluxo $f$ é dito viável se satisfaz duas condições (veja~\cite{ford1956maximal, cook1998combinatorial}):

\begin{enumerate}
    \item \textbf{Restrição de Capacidade:} O fluxo em cada arco é não negativo e não excede sua capacidade:
    \begin{equation}
        0 \le f(u,v) \le c(u,v), \quad \forall (u,v) \in A
    \end{equation}

    \item \textbf{Conservação de Fluxo:} Para todo vértice intermediário $v \in V \setminus \{s, t\}$, a soma dos fluxos que entram em $v$ é igual à soma dos fluxos que saem de $v$:
    \begin{equation}
        \sum_{(u,v) \in A} f(u,v) = \sum_{(v,w) \in A} f(v,w)
    \end{equation}
\end{enumerate}

O valor do fluxo $f$, denotado por $|f|$, é o fluxo líquido que sai da fonte $s$:
\begin{equation}
    |f| = \sum_{(s,v) \in A} f(s,v) - \sum_{(v,s) \in A} f(v,s)
\end{equation}
Pela conservação de fluxo, este valor é igual ao fluxo líquido que entra no sumidouro $t$. O Problema do Fluxo Máximo consiste em determinar um fluxo viável $f$ que maximize $|f|$.

Essa estrutura é formalizada computacionalmente na classe base \code{FlowNetwork}, que gerencia a estrutura de adjacência da rede e assegura que as restrições de capacidade e conservação sejam mantidas.


